\documentclass{presentation}

\title{Άλγεβρα Α Λυκείου}
\subtitle{Προτάσεις και λογικοί σύνδεσμοι}
\author[Λόλας]{Κωνσταντίνος Λόλας}
\institute[$10^ο$ ΓΕΛ]{$10^ο$ ΓΕΛ Θεσσαλονίκης}

\begin{document}

\begin{frame}
  \titlepage
\end{frame}

\section{Θεωρία}

\begin{frame}{Η Άλγεβρα κάνει δηλώσεις}
  \begin{itemize}[<+->]
    \item «Το $2+3$ κάνει $5$.»
    \item «Αν $x=4$, τότε $x^2=16$.»
    \item «Το $x$ είναι μεγαλύτερο από $2$ και μικρότερο από $5$.»
  \end{itemize}

  \begin{exampleblock}{Σήμερα}
    Θα μάθουμε να διαβάζουμε αυτές τις δηλώσεις χωρίς να χρειάζεται
    αποκωδικοποιητής.
  \end{exampleblock}
\end{frame}

\begin{frame}{1. Πρόταση}
  \begin{block}{Πρόταση}
    Είναι μια πρόταση που είναι είτε \textbf{αληθής} είτε \textbf{ψευδής}.
  \end{block}

  \begin{itemize}[<+->]
    \item $2+3=5$ \hfill Αληθής
    \item $10<4$ \hfill Ψευδής
    \item «Λύσε την άσκηση!» \hfill Δεν είναι πρόταση
    \item $x+1>3$ χωρίς γνωστό $x$ \hfill Δεν αποφασίζεται ακόμη
  \end{itemize}

  \begin{alertblock}{Γρήγορος έλεγχος}
    Μπορώ να απαντήσω «Αληθές» ή «Ψευδές»; Τότε έχω πρόταση.
  \end{alertblock}
\end{frame}

\begin{frame}{Πρόταση με γράμμα}
  Ονομάζουμε τις προτάσεις με κεφαλαία γράμματα:
  \[
    A: x=3, \qquad B: x+2=5.
  \]

  Για $x=3$ ισχύουν και οι δύο. Για $x=1$ δεν ισχύει καμία.

  \begin{exampleblock}{Μικρή προσοχή}
    Το $A$ είναι ολόκληρη η πρόταση. Το $x$ είναι απλώς ο άγνωστος.
    Μην τα γνωρίσετε και μπερδευτούν!
  \end{exampleblock}
\end{frame}

\begin{frame}{2. Συνεπαγωγή}
  \begin{block}{«Αν $A$, τότε $B$»}
    \[
      A\Rightarrow B
    \]
  \end{block}

  \begin{exampleblock}{Παράδειγμα}
    \[
      x=4\Rightarrow x^2=16
    \]
    Το $x=4$ αρκεί για να πάρουμε $x^2=16$.
  \end{exampleblock}

  \begin{alertblock}{Το βέλος έχει φορά!}
    \[
      x^2=16\not\Rightarrow x=4
    \]
    Μπορεί να είναι και $x=-4$. Το βέλος δεν κάνει αναστροφή χωρίς φλας.
  \end{alertblock}
\end{frame}

\begin{frame}{3. Ισοδυναμία}
  \begin{block}{«$A$ αν και μόνο αν $B$»}
    \[
      A\Leftrightarrow B
    \]
    Δηλαδή ισχύουν και οι δύο κατευθύνσεις:
    \[
      A\Rightarrow B \quad\text{και}\quad B\Rightarrow A.
    \]
  \end{block}

  \begin{exampleblock}{Παράδειγμα}
    \[
      x^2=0\Leftrightarrow x=0.
    \]
    Οι δύο πλευρές έχουν ακριβώς τις ίδιες λύσεις.
  \end{exampleblock}
\end{frame}

\begin{frame}{4. Το «ή»}
  \begin{block}{Διάζευξη}
    \[
      A\vee B
    \]
    σημαίνει «$A$ ή $B$». Αρκεί να ισχύει τουλάχιστον μία.
  \end{block}

  \begin{exampleblock}{Παράδειγμα}
    \[
      x<0\ \text{ή}\ x=0
      \quad\Longleftrightarrow\quad x\le0.
    \]
    Στα μαθηματικά το «ή» επιτρέπει και τις δύο περιπτώσεις. Γενναιόδωρο!
  \end{exampleblock}
\end{frame}

\begin{frame}{5. Το «και»}
  \begin{block}{Σύζευξη}
    \[
      A\wedge B
    \]
    σημαίνει «$A$ και $B$». Πρέπει να ισχύουν και τα δύο.
  \end{block}

  \begin{exampleblock}{Παράδειγμα}
    \[
      x\ge2\ \text{και}\ x\le5
      \quad\Longleftrightarrow\quad 2\le x\le5.
    \]
    Το $x$ περνάει από δύο ελέγχους. Δεν αρκεί να περάσει μόνο τον έναν!
  \end{exampleblock}
\end{frame}

\begin{frame}{6. Άρνηση}
  \begin{block}{Το «όχι»}
    Η άρνηση της $A$ γράφεται $\overline{A}$ και κάνει την αληθή πρόταση ψευδή
    και την ψευδή αληθή.
  \end{block}

  \begin{exampleblock}{Στις ανισώσεις}
    \[
      \overline{(x>3)}\Leftrightarrow x\le3,
      \qquad
      \overline{(x\ge2)}\Leftrightarrow x<2.
    \]
  \end{exampleblock}

  \begin{alertblock}{Δεν αλλάζουμε μόνο το σύμβολο}
    Η ισότητα πηγαίνει κι αυτή στην απέναντι πλευρά.
  \end{alertblock}
\end{frame}

\begin{frame}{7. Αντιθετοαντιστροφή}
  \begin{block}{Ο κανόνας}
    \[
      A\Rightarrow B
      \quad\Longleftrightarrow\quad
      \overline{B}\Rightarrow\overline{A}.
    \]
  \end{block}

  \begin{exampleblock}{Παράδειγμα}
    \[
      x>5\Rightarrow x>2
    \]
    έχει αντιθετοαντίστροφη:
    \[
      x\le2\Rightarrow x\le5.
    \]
  \end{exampleblock}

  \begin{alertblock}{Η συνταγή}
    Αρνούμαι το $B$, αρνούμαι το $A$, αλλάζω σειρά. Έτοιμο!
  \end{alertblock}
\end{frame}

\begin{frame}{Η κάρτα-σκονάκι}
  \begin{center}
    \begin{tabular}{c|c}
      Λόγια & Σύμβολο\\ \hline
      αν $A$, τότε $B$ & $A\Rightarrow B$\\
      $A$ αν και μόνο αν $B$ & $A\Leftrightarrow B$\\
      $A$ ή $B$ & $A\vee B$\\
      $A$ και $B$ & $A\wedge B$\\
      όχι $A$ & $\overline{A}$
    \end{tabular}
  \end{center}

  \begin{block}{Σκέψου πρώτα τα λόγια}
    Μετά διάλεξε το σύμβολο. Τα σύμβολα είναι βοηθοί, όχι εξεταστές!
  \end{block}
\end{frame}

\moodle

\section{Ασκήσεις}

\exercises

\begin{askisi}
  \textbf{Το διάλειμμα του $x$.} Έστω $x\in\mathbb{R}$ και
  \[
    A: x\ge2,\qquad B: x\le5,\qquad C: x=4.
  \]

  \begin{enumerate}[<+->]
    \item Ποιο είναι το σύνολο λύσεων της $A\wedge B$;
    \item Τι σημαίνει η $\overline{A}\vee\overline{B}$;
    \item Είναι σωστή η $A\wedge B\Rightarrow C$; Βρες ένα αντιπαράδειγμα.
    \item Γράψε την αντιθετοαντίστροφη της προηγούμενης συνεπαγωγής.
    \item Δείξε ότι $C\Leftrightarrow(x\ge4\wedge x\le4)$.
  \end{enumerate}
\end{askisi}

\begin{lisi}
  \textbf{Συζήτηση}
  \begin{enumerate}[<+->]
    \item $A\wedge B$ σημαίνει $2\le x\le5$, άρα $x\in[2,5]$.
    \item $\overline{A}\vee\overline{B}$ σημαίνει $x<2$ ή $x>5$.
    \item Όχι. Για $x=3$ ισχύει $A\wedge B$, αλλά όχι $C$.
    \item Η αντιθετοαντίστροφη είναι
          $\overline{C}\Rightarrow\overline{(A\wedge B)}$, δηλαδή
          $x\ne4\Rightarrow(x<2\vee x>5)$.
  \end{enumerate}
\end{lisi}

\begin{frame}{Το τελευταίο βήμα}
  \[
    x=4\Leftrightarrow(x\ge4\wedge x\le4).
  \]

  Για να ισχύουν και οι δύο ανισώσεις, το $x$ δεν έχει πού αλλού να πάει:
  \[
    x=4.
  \]

  \begin{block}{Συμπέρασμα}
    Το «και» κρατά το κοινό μέρος, το «ή» ενώνει περιπτώσεις και η
    αντιθετοαντιστροφή γυρίζει τον συλλογισμό χωρίς να τον χαλάει.
  \end{block}
\end{frame}

\end{document}
